\chapter{Einleitung}

Um nicht sichtbare und nicht geladene radioaktive $\gamma$-Strahlung zu untersuchen, die aufgrund ihrer großen Frequenz und daher ihrer hohen Energie zu dauerhaften biologischen Schäden bei Exposition führen kann, werden Detektoren auf Basis von Szintillation in einem dotierten Kristall oder Ladungsansammlung in einem Halbleiter genutzt. Hierbei wird die Wechselwirkung der ungeladenen $\gamma$-Photonen mit Materie in Form von Photoeffekt, Coulomb-Streuung und Paarbildung ausgenutzt, wobei die so an das Material abgegebene Energie gemessen werden kann. Aus der so bestimmten Energie kann schließlich auf die Radionuklide geschlossen werden, deren Zerfälle jeweils mit charakteristischen $\gamma$-Photon-Emissionen verbundenen sind \cite[518-520]{gammahandbook}.

Im durchgeführten Versuch sollen mit einem \textit{Natrium-Jodid}-Szintillationsdetektor (NaI) sowie einem \textit{High-Purity Germanium}-Halbleiterdetektor (HPGe) jeweils die Energiespektren der emittierten Photonen innerhalb des $\gamma$-Strahlungsfrequenzbereich über $0$ bis $\SI{1550}{\kilo\eV}$ von verschiedenen Proben aufgenommen werden \cite[1]{521versuchsanleitung}.
Als Proben werden $\ce{ ^60Co }, \ce{ ^137Cs }, \ce{ ^152Eu}$ \cite[3]{521versuchsanleitung} und eine unbekannte Bodenprobe (diese nur mit dem HPGe-Detektor als Langzeitmessung über 18 Stunden) vermessen sowie jeweils Untergrundmessungen der gemessenen Signale der Detektoren ohne Probe durchgeführt. Aus den Messungen sollen die Energie-Auflösungen der spezifischen verwendeten Detektoren,  ihre \textit{peak-to-total}-Verhältnisse sowie ihre energieabhängige Effizienzen bestimmt werden.
Außerdem soll mithilfe des Spektrums der Bodenprobe bestimmt werden, welche Radionuklide in ihr vorhanden sind.
